% Generated draft: requires mathematical and source review.
\documentclass[11pt]{article}
\usepackage[margin=25mm]{geometry}
\usepackage{fontspec}
\setmainfont{Latin Modern Roman}
\usepackage{amsmath,amssymb,mathtools,unicode-math}
\setmathfont{Latin Modern Math}
\usepackage{enumitem,needspace,xurl,hyperref}
\setlength{\parindent}{0pt}
\setlength{\parskip}{5pt}
\setlength{\emergencystretch}{3em}
\newcommand{\N}{\mathbb{N}}
\newcommand{\Z}{\mathbb{Z}}
\newcommand{\Q}{\mathbb{Q}}
\newcommand{\R}{\mathbb{R}}
\newcommand{\C}{\mathbb{C}}
\newcommand{\F}{\mathbb{F}}
\newcommand{\A}{\mathbb{A}}
\newcommand{\PP}{\mathbb{P}}
\newcommand{\E}{\mathbb{E}}
\newcommand{\eps}{\varepsilon}
\newcommand{\dd}{\,\mathrm{d}}
\begin{document}
\section*{UnsolvedMath ID: 30006393}
\textbf{Source problem code:} catalogue entry
\subsection{Definitions and mathematical statement}
Let \(p\) be a prime, and let \(\F_p=\Z/p\Z\) be the field with \(p\) elements, with its additive group structure. For a subset \(S\subseteq\F_p\setminus\{0\}\), write \(k=|S|\). An ordering of \(S\) is a sequence \((s_1,\ldots,s_k)\) containing each element of \(S\) exactly once. Its nonempty successive partial sums are
\[
 t_i=\sum_{r=1}^{i}s_r\in\F_p,\qquad 1\leq i\leq k.
\]
The question in the supplied cleaned statement is whether
\[
 \forall p\text{ prime}\;\forall S\subseteq\F_p\setminus\{0\}\;\exists(s_1,\ldots,s_k)\text{ ordering }S\;\forall\,1\leq i<j\leq k,\quad t_i\ne t_j.
\]
For \(S=\varnothing\), the empty ordering satisfies this condition vacuously. No restriction on \(\sum_{s\in S}s\) is supplied.

\emph{Partial-sum convention requiring review.} This formulation interprets successive partial sums as \(t_1,\ldots,t_k\), without including the empty partial sum \(t_0=0\). The supplied record does not explicitly settle that convention. Requiring \(t_0,t_1,\ldots,t_k\) to be pairwise distinct would change the statement and would be incompatible with any nonempty allowed set whose total sum is zero.

\emph{Scope and supplied status.} The extracted question concerns additive prime fields. The original passage also discusses a possible extension to every finite group, but does not supply the preceding question, its group-operation conventions, the statement of Theorem 1, or the bounds defining the intermediate regime. Those broader references do not determine an additional precise conjecture here. The passage's assertion that the intermediate regime remains open is qualified by the supplied dated assessment: as of its stated check date, 2026-08-22, the curation attributes a resolution for all sufficiently large primes to Pham and Sauermann, combining small- and large-subset results, while describing the broader arbitrary finite-group question as open. These are supplied literature assessments, not independently verified mathematical conclusions.
\subsection{Short English statement}
For each prime \(p\), can every subset of the nonzero elements of \(\F_p\) be ordered so that its nonempty successive partial sums are pairwise distinct? This draft excludes the empty partial sum; that convention needs source review. The supplied curation labels the problem partially solved and reports a result for sufficiently large primes, with a dated assessment of 2026-08-22.
\subsection{Sources}
\begin{itemize}
\item Supplied local problem registry: permanent UnsolvedMath ID \(30006393\), title \emph{Graham's Rearrangement Conjecture over Prime Fields}, exact source problem code \texttt{catalogue entry}. The record supplies the status \texttt{partially\_solved} and statement status \texttt{corrected\_verified}. Its verification field reports that the official report was checked and the extracted notation repaired; this draft does not independently verify that report or certification.
\item Supplied report citation: \emph{Combinatorics, Probability and Computing (2026)}, Oberwolfach Reports, DOI 10.4171/owr/2025/42, pp. 2243--2322. \url{https://doi.org/10.4171/owr/2025/42}. The record separately supplies proposed year 2025 and an \texttt{owr-year:2025} tag.
\item Historical attribution in the supplied original passage: the problem was first posed by Graham in 1971, with internal reference label [10], and later reiterated in an open-problems book of Erd\H{o}s and Graham, with internal label [8]. The passage also mentions Gordon's setting, a directed variant of an Andersen conjecture with internal label [2], and Theorem 1 of the surrounding report. Complete bibliographic entries and the surrounding formulations are not supplied.
\item Supplied literature source: \emph{On Graham's Rearrangement Conjecture}. \url{https://arxiv.org/abs/2602.15797}. The supplied assessment attributes the sufficiently-large-prime resolution to Pham and Sauermann.
\item Supplied literature source: \emph{On Graham's Rearrangement Conjecture over \(\F_2^n\)}. \url{https://arxiv.org/abs/2508.18254}.
\item Literature assessment checked at \texttt{2026-08-22}: supplied dated curation, not independently verified. It labels the problem partially solved, reports that the intermediate-size regime for prime cyclic groups is no longer wholly open because of the sufficiently-large-prime result, and describes the broader arbitrary finite-group question as open. The background explicitly requests independent verification of these time-sensitive claims.
\end{itemize}
\end{document}
